\quellseite{487}
\chapter[§ 34. THE METHOD IN BIBLICAL HISTORY]{§ 34. THE METHOD OF INSTRUCTION IN BIBLICAL HISTORY}

Compare the literature in § 25 and 32. A. Eckert, pp. 79–118. J. Gottschick, pp. 175–181. O. Harnisch, pp. 63–64. *W. Bittorf, \emph{Methodik}, ²1908, pp. 32–118. J. Berndt, pp. 78–95. R. Kabisch, pp. 182–219. J. Steinbeck, pp. 185–193. *F. W. Doerpfeld, \emph{Ges. Schriften iii: Der Religionsunterricht}, ²1895. P. Staude, \emph{Die formalen Stufen i. bibl. Geschichtsunterricht} in Kehr’s Paedag. Blaetter, 1884. L. Hohmann \emph{Methodik d. ev. Religionsunterrichts}, ²1904. *F. Gansberg, \emph{Schaffensfreude}, ³1909. *H. Scharrelmann, \emph{Weg zur Kraft}, ⁷1910. *O. and E. Zurhellen, \emph{Wie erzaehlen wir den Kindern die biblischen Geschichten}, ²1910. Max Paul, \emph{Fuer Herz und Gemuet der Kleinen}, ⁶1911. M. Reu, \emph{Grundsaetze z. Herstellung von Sonntagsschulliteratur} in Kirchl. Ztschr., 1911. *M. Reu, \emph{How I Tell the Bible Stories to My Sunday School}, 2 vols., ²1926. \textsc{Teaching Helps for All Grades:} K. Francke, \emph{Weide meine Laemmer. Die Heilige Geschichte der Jugend erzaehlt}, 1897. O. Zuck, \emph{Der gesamte Religionsunterricht}, 3 parts, 5 vols., ⁵1902 ff. R. Staude, \emph{Praeparationen zu der bibl. Gesch. d. Alten u. Neuen Tests.}, 3 vols., ⁴1905 ff. E. Thraendorf and H. Meltzer, \emph{Der Religionsunterricht}, 3 parts, ²1905 ff. *Reukauf and Heyn, \emph{Evang. Religionsunterricht}, vol. iii to ix, ⁴1907 ff. Kessel and Spanuth, \emph{Praepar. f. d. ev. Religionsunterricht}, 3 parts, ²1910 f. *G. Fankhauser, \emph{Die bibl. Gesch. in Sonntagsschule u. Religionsstunde}, ³1915. G. Fankhauser, \emph{Christ d. Retter. Gesch. aus d. Leben Jesu. kleinen u. grossen Kindern erzaehlt}, 1925. *G. Fankhauser, \emph{Der Weg zum Kinde. Monatsschrift f. christl. Erziehung u. bibl. Unterricht}, 1915 ff. \textsc{Teaching Helps for the Instruction of Beginners:} G. L. W. Koehnlein, \emph{Bibl. Geschichten f. Kinder von 4 bis 9 Jahren}, 1854. J. F. Ranke, \emph{Des Kindes erster Unterricht aus Gottes Wort angeschlossen an die 40 Bilder der Kaiserswerther Bilderbibel. Eine Vorstufe z. d. ersten Religionsunterricht}, ²1873. F. Wiedemann, \emph{Wie ich meinem Kleinen die bibl. Geschichten erzaehle}, about 1865. Graefin Poninska, \emph{Biblische Geschichten. Nach d. Schnorrschen Bilderbibel erzaehlt f. kleine Kinder}, 1879. L. Wangemann, \emph{Der erste biblische Anschauungsunterricht. Anweisung z. Gebrauch der 20 Anschauungsbilder}, ²1907. P. Staude, \emph{Praepar. f. d. ersten Religionsunterricht}, 1889 f. A. Gieseler, \emph{Der Religionsunterricht auf. d. Unterstufe}, 1900. G. Kaelker, \emph{Bibl. Anschauungsunterricht},
\quellseite{488}
²1905. E. Zeissig, \emph{Von d. bibl. Gesch. auf d. Unterstufe in Katech. Zeitschr.}, 1907. R. H. Evers, \emph{Die bibl. Geschichte u. deren Behandlg. auf d. Unterstufe}, ³1909. M. R. Unger, \emph{Beitrag zur erbaulichen Behandlg. des Religionsunterrichts auf der Unterstufe der Volksschule}, ab. 1912. *M. Reu, \emph{Fuer die Kleinen}, 1913. K. Schneeberg, \emph{Uns’ Herrgott un sin Lued}, 1914. Mrs. H. Gaskoin, \emph{Children’s Treasury of Bible Stories}, 1896. James Baldwin, \emph{Old Stories of the East}, 1896. R. Bird, \emph{Jesus the Carpenter and Joseph the Dreamer}, 1901. Laura E. Ragin, \emph{Kindergarten Bible Stories}, 1905. Geo. Hodges, \emph{The Garden of Eden, The Castle of Zion, and When the King Came}, 1904 ff. R. G. Moulton, \emph{The Art of Telling Bible Stories}, 1904. *Sara Coni Bryant, \emph{How to Tell Stories to Children}, 1905. Rich. T. Wyche, \emph{Some Great Stories and How to Tell Them}, 1910. *E. P. St. John, \emph{Stories and Story-Telling}, 1910. M. Reu, \emph{For Beginners. Fifty Stories for the Little Ones}, 1916. M. Affolter, \emph{Kommt, Kindlein, hoert}, ²1926. Leipziger Lehrerverein, \emph{Im Strome des Lebens, Sammlung von Erzaehlungen und Geschichten zur Belebung der religioesen Jugendunterweisung}. Sneath, Hodges, and Tweedy, \emph{The King’s Highway Series}, now 8 vols., containing much valuable material for telling the Bible stories and for the application upon the life of the child.

That the material assigned for any particular lesson in Biblical History should not be loose conglomeration of isolated events, but a catechetical unit, has already been shown (p. 424). If this statement is expressly repeated here, it is done for the reason that the truth involved is the postulate for all that is to be said in this paragraph concerning the method underlying instruction in Biblical History. While it is quite permissible in the upper grades or on the occasion of a review to place a number of narratives under a common point of view and thus to weave them into a unit, the catechetical unit for the first treatment is always the single narrative. How to treat the single narrative, then, will be the chief problem. Many are of the opinion that the main object of instruction in Biblical History is merely the memorizing of its principal contents, for the one purpose of having an opportunity to refer, in sermon and catechesis, to the material stored up in the memory. If that view were correct,
\quellseite{489}
the science of catechetics might well refrain from investigating the various methods as to their appropriateness. In that case we might be satisfied with reading the story several times and conducting a quiz for the purpose of ascertaining whether there has been an appropriation of the contents. But if instruction in Biblical History is to do its share toward reaching the aim of religious instruction and training as it is given on page 280; if it is to be of present value for the children, if it is to influence their whole inner life and to implant the desire to live according to God’s will, if it is to become a dynamic factor in their young lives, then the question as to the didactic treatment of the individual biblical narrative proves of uncommon importance, and requires specific treatment.

In order to attain this purpose, there will have to be in evidence what we described in §32 as the essential features of the whole process of instruction. In other words, the treatment of the individual narrative should cover the three steps of presentation, penetration, and application. The penetration will also comprise what the Herbartians call “association” (cf. p. 448) in case the latter seems pedagogically valuable; and invariably it will issue in the “generalization” of the specific truth educed from the individual narrative (fourth step of the disciples of Herbart). Again, presentation should be preceded by a statement of aim and frequently also by a brief preparation.

\section{1. Statement of aim}

A sermon that pursues only the general aim of edifying the congregation, will degenerate into an aimless and useless talk; every time the preacher ascends the pulpit he must have a clear and specific aim, and he will say nothing except what serves to attain this aim. His success is proportionate to the clearness with which he formulates his aim and to the “goalwardness” of the sermon. Whether he states his aim in the form of a theme or not—he must
\quellseite{490}
have one and constantly pursue it. The same is true of the catechist. He should never present a Bible story to his children without determining exactly what he means to accomplish. But should he specify and name his aim at the very beginning of the catechesis? Certainly not the inner aim, the ultimate aim to be realized with his pupils. In that case he would anticipate the best that is to come and himself thwart the children’s attention; he would hand over to them what they should have gained by earnest efforts, and he would deprive himself of the best opportunity to engage their mental co-operation. What he will do, however, is to state the aim formally. He will not say: “Today let us learn that nothing so pleased God in Abraham as faith”, but: “Let us learn today what pleased God most in Abraham.” The former is the announcement of the inner aim, the latter, of the outward one; the former is the material announcement of the aim, the latter, the formal. The latter is necessary in order to concentrate and focus the scattered thoughts of the children, and to direct them to the point on which the whole catechesis is to turn. The expediency of stating the aim must be recognized by anyone who realizes the variety of thoughts with which the children are still occupied when the Bible History period begins and who remembers that the “narrowness of consciousness” (cf. p. 207) prevents the soul from thinking of more than one subject at a time. As long as the soul clings in thought to the old, I cannot fill it with anything new. It will behoove me first to detach it from the old and lead it in the direction of the new: this is the very purpose of the statement of aim. In proportion as the catechist succeeds in making the statement attractive, brisk, interesting, he will accomplish his purpose. Statements of this kind would be: How sin came into the world; How God made man happy; Of Joseph, Jacob’s favorite son; How Joseph came to Egypt; How a shepherd lad
\quellseite{491}
was elected king; How little David fought with the tall giant; Of David’s dearest friend Jonathan; How the lad Jesus came into the temple for the second time; Of a man who fell among robbers; Jesus takes care even of the lost; Even believers may yet perish; How the first black man became a Christian, etc. Let the student compare the headings of the Wartburg Lesson Helps (Second Course; but note that the inner, or material aim has in several instances purposely been used there as heading). In connection with the statement of aim it should also be remembered that the child, and the common man in general, prefer the use of verbs to that of nouns. “The creation of the world” does not mean so much to them as “How God created the world.” Sometimes the statement of aim of the new story is readily suggested by the preceding one. For instance, the children’s sympathy has been aroused if the teacher has, with some degree of success, told the story how Joseph was sold into slavery; they are genuinely interested in his welfare; and on their lips lies the question: What became of Joseph in Egypt? Could the teacher desire a clearer statement of aim for the following story?

\section{2. Preparation}

It is often advisable that the presentation is preceded by a preparation. In the statement of aim a word has perhaps been used which points the child to conditions in life with which he is well acquainted. The discussion of these will supply the apperceptive faculty with the fulcrum which is to bear the weight of the new material in the narrative to be treated (p. 212). Should, e. g., the aim of the narrative have been stated as: “How the first black man became a Christian,” it is advisable to remind the child of the negro whom he has seen on the street, possibly also of the fact that even negroes can be true Christians. If, in addition, a short story should be told of some negro known as an earnest Christian, the child has not only been
\quellseite{492}
provided with the material which enables him to understand the statement of aim, but a torch has been lit which shines into the distance so that the erstwhile nebulous and obscure form of the Ethiopian eunuch has become something concrete and discernible. It is also possible that the new narrative deals with objects and affairs altogether remote from the present viewpoint and concept sphere of the children, but, nevertheless, of fundamental moment for an intelligent comprehension of that particular narrative. Such matters cannot be explained incidentally during the process of presentation; for this reason they constitute material for the preparation. As a case in point, in the story of the anointing of David for the royal office the concept “anointing” occupies a prominent place. The catechist will find far fewer obstacles to a successful disclosure of the thought underlying this narrative and the development of the inner aim (namely 1 Sam. 16: 7), if, in the process of presentation and penetration, he will not be compelled to dwell upon an explanation of the custom of anointing. In the preliminary discussion he will speak of the inauguration of our presidents, of the crowning of emperors, or of the practice of anointing priests and kings in Israel and the purpose and meaning back of it. Or, the catechist, very properly, may feel justified in connecting the new narrative with the one just disposed of, because only thus the children can be put into the right frame of mind to bring the needed amount of interest to bear upon the new material. It is the province of the preparation to bring about such a frame of mind. If, for instance, the sale of Joseph into captivity has been the subject of the previous lesson, the children should be reminded anew how Joseph must have felt on his way to Egypt—leaving his father, running beside the camel in the hot sand of the desert, sobbing and moaning in his lonesome grief, telling himself that he would probably never see his home again. Now the
\quellseite{493}
children are mentally prepared to follow Joseph’s further experiences in Egypt. Thus the preparation aims to enhance the interest of the child in the new story already aroused by the statement of aim; it will remove all obstacles in the way of an effectual appropriation of the new material, and attune the cords of the soul so that the new impressions may play upon it with power and without interference. It will often be well to let the preparation issue in a re-statement of the aim. — So far as the outward form is concerned, the preparation is not necessarily purely acroamatic; it is more likely to bear the character of an informal conversation. While the explanatory words of the catechist will usually form the bulk of the preparation, there may be instances when, to make it effectual, the whole preparation will consist of questions and answers (cf. the “method of \emph{construction}” discussed below). The “preparation” may be omitted entirely when the statement of aim and the subsequent presentation are immediately intelligible; at any rate, this is better than artificially to drag it in. — Compare the practical examples appended to the fifth part of this textbook.

\section{3. Presentation}

How the new narrative is presented, is of the greatest importance for the realization of the aim of instruction or training. If the presentation fails to arouse interest, to stir the emotions, to move the will, the other steps of the teaching process are not likely to have better effects, and the instruction as a whole remains without practical results. There are three modes of presenting the new material, each one having its advocates among the ranks of the educators. The new story may be offered through reading, narration, or “construction.”

\emph{Reading} is the most ineffectual form imaginable, especially when the children are expected to read the story themselves from the textbook or the Bible. Even though they be able to read with fluency and expression, it is impossible
\quellseite{494}
for them to receive so clear an intuition as when the story is read to them in an exemplary manner; and reading the Bible narrative over once, in most cases, will have no influence upon their inner life. But does it follow from this that the teacher is to read the story to the children? It must be conceded that, if the teacher reads well, the interest of the children can be aroused. But equally correct, and even more important, is what Schueren says upon this point (p. 15): “Whatever the teacher exacts from the children, he must first exact from himself. How repugnant it is when the teacher says to the children: “You should know this, even though I myself do not.” But that is exactly what he implies, when he reads the story to the children. He should not even use reference notes, since he does not allow the children to use these when the subject is reviewed. Even more injurious than setting this example of indolence is the resulting lack of discipline and the loss of animation in teaching. Children should know themselves under constant observation, if order is to be maintained among them. When the teacher reads to the children, he is unable to watch them, unable to nip disturbance in the bud;—and disturbance is likely to be particularly frequent in just this case, for the reason that a lifeless instruction possesses no fascination for the children. If there is to be life in instruction, the subject matter must live in the teacher instead of merely being in the book. How can the teacher arouse and maintain in the children the conviction that he himself believes the material that he presents, if he resorts to reading? Many other things besides the voice should speak during instruction, notably the eye. There is no telling what effect may be produced by a soulful eye turned to the hearer. The eye of the teacher should rest upon that of the child, and the eye of the child upon that of the teacher. Often the lips will strive in vain to utter what the eyes express without
\quellseite{495}
effort.” Furthermore, what text should the teacher select for reading? That of Scripture? This is the view of many, who think that veneration for Holy Scripture does not permit any choice whatever. But the Scriptures have not been written for the children in the sense that their language and mode of presentation is always familiar to them. On the contrary, the language of the Scriptures often requires to be adjusted to the needs of children, and to be translated into their own familiar speech. In order to bring home to the children the actual sense, it will occasionally be necessary to leave out parts of the Scripture text, and, again, to supplement it by explanatory comments. In due recognition of this fact, Bible Histories have been compiled with texts frequently quite different from that of the Scriptures, so far as the arrangement of words is concerned, texts that rise from comparative freedom from the Scriptural text in the lower grades to increasing conformity to it in the upper grades (compare the three grades of the Second Course of the Wartburg Lesson Helps). This practice is warranted without doubt; and it is even necessary in view of the needed reviews and drills; nevertheless, the objections to presentation by reading stated above apply also here; moreover, the mere reading of the story from the textbook does not suffice to unfold completely before the pupils the latent powers of the text of Scripture and to bring them to bear upon the threefold life of the soul.

The reason is this: In every story, the Bible story not excepted, a fourfold factor is to be discerned: (1) the historic act as such; (2) the framework of the story—for instance, certain features pertaining to the respective fields of geography, natural history, and cultural history within which the action moves; (3) the inner world of thoughts and motives of the story, i.e., the thoughts, emotions, reflections, motives, etc., actuating the persons described; and (4) the
\quellseite{496}
religious and moral nature of such thoughts and motives, which endows the action with its real value. These four factors must be comprehended by my pupils; until that has been achieved, I have not fully presented the story nor utilized all of its dormant powers. The ordinary school texts, however, permit only the first one to come into its own; the second falls short and has to be explained afterward whereby the effect of the story is marred; the third and the fourth are merely suggested; and so far as the fourth is concerned, it often is totally ignored. And yet, the educative force of the story is largely found in the last two factors. The pupil must be trained not to stop at externals in judging of an action, but to make his way to the motive of which it is an expression, and thus to form a judgment concerning its religious and moral value (pp. 223 f., 279 f.). To this must be added that the schooltext as well as that of the Bible often indicates important situations with a single stroke. This may be, and as a matter of fact often is, decidedly characteristic and significant; but the child does not understand it in its bearings, and thus loses something of great value. For instance, what anguish and despair, what remorse and futile hope, what horror of death and silence of the grave are expressed in that brief word of Gen. 7:21. 22! All this should be brought out; the class must vividly be made conscious of it. Or consider the situation sketched by Acts 9:22: “Paul confounded the Jews that dwelt at Damascus and proved that this is the Christ”; these words are only meaningless sounds for the class, probably not even understood in their literal sense, until the teacher unfolds them and presents to his pupils a vision of Paul’s evangelistic activity with its overwhelming arguments from the sphere of prophecy and fulfillment. Facts such as these call for a mode of presentation different from a mere reading of the respective text by the pupils or by the teacher. Even
\quellseite{497}
when subsequently the penetration supplies many things neglected by the presentation, this is a poor makeshift: until that time comes, the child has to worry along with much undigested material; no overwhelming effect upon his inner life, particularly his will and emotions, has been produced. Whether the soul will later be equally susceptible to such influence, is questionable: the first impression is usually also here decisive.

For the reasons given, many catechists today, instead of reading the text or having it read, prefer the method of \emph{construction}. They understand by this the method through which the teacher, in the form of a conversation, introduces the children to the situation postulated by the story, and then, by questions and answers, prompts them to find the several features of the story and thus in a sense to construct it for themselves. Where this method obtains, the above mentioned four characteristic features of the story, which require due attention if the story is to fulfill its purpose—the last three in particular—come into their own. Notwithstanding, this treatment of the story is still inadequate. While the mental activity of the children could hardly find a better field for independent exertion, the biblical story cannot, where this method is followed, produce the designed effect upon the emotional and volitional life of the children; and when climaxes are reached in the history of revelation and God’s redemptive acts are the heart of the story, this method fails completely. How should children discover God’s purpose in regard to man either in the past or present? Even though the thesis should be given to them: “God is love”, all attempts to discover His ways from that premise would result in aimless groping, but never in a knowledge of the ways which God has actually taken. If, however, this dangerous deviation is avoided and on all important points information about God’s redemptive work is imparted to
\quellseite{498}
the children, then it is thereby tacitly admitted that a story hardly requires any other form of presentation than that of narration. While the constructive method has originated in an exaggeration of the principle of training the child for mental self-activity, it is not without its element of truth, in that the catechist, who presents the subject matter in the form of narration, will find it expedient now and then to interrupt his narration by a question, which is intended not only to make sure of the pupils’ attention, but also to enable them to carry the action forward themselves or to give an account of the thoughts and motives underlying it.

But \emph{how should the story be told?} Here, too, a three-fold mode is advocated. (a) There are those who think that it ought to be told in precisely the same manner as the pupils shall afterward find it in the school-text, with only this exception that brief explanatory remarks may be interspersed. That this view is hardly correct, is readily inferred from what has been said above; for most of the arguments set forth against the reading of the Bible story by the teacher, apply to this method of narration as well. (b) Others have advocated a method diametrically opposite, in that they demand a narration that is not only altogether independent of the Bible text but also modernized to such a degree that the very coloring of the Bible is discarded, and the Bible characters appear like present-day country or city folks. While traces of this method are found already in Paul and Zurhellen, it is Gansberg and Scharrelmann who have consistently carried it out. Not only “full detailedness and exhaustive motivation,” but also “the most child-like modernization” is Scharrelmann’s motto. He says: “When I am unable to depict ancient Jewish and Egyptian conditions and scenes, I substitute modern ones for them,” and he follows this principle even more consistently than the great medieval painters by whose methods he justifies his own. We do not
\quellseite{499}
overlook the fact that these demands contain elements of truth which we ourselves recognize and adopt (v. i.); nevertheless, we must repudiate this mode of narration or presentation; for, to us, not only the soul of the child in whose behalf the demand is made is holy, but the material as well. Much as we emphasize the necessity of setting forth the truths of Christian faith and life contained in the several Bible narratives in such a way that the child’s interest is aroused on every hand, we are at the same time aware of the Bible story’s title to reality. It is to us an educational factor for the very reason that it is a link in the chain of that history which prepares or provides our salvation; if it were a mere vestment for garbing moral thoughts, we should not use it at all as an educational factor and prefer to have recourse to tales or stories from the present.

The proper mode of presentation is (c) \emph{free narration}. In far-reaching independence of the Biblical wording such narration gives a clear, vivid, and exhaustive picture of the event; it gives due attention to the progress of the action and to the whole “frame work of the story”; and tenderly it discloses the inner world of thoughts and motives and their moral and religious bearings. At the same time, however, in spite of all independence from the Biblical wording, in spite of all detailedness and motivation, it is in point of content, faithful to Scripture and in point of form noble and chaste; it intends to set in motion the intellectual, emotional, and volitional soul-life of the child by no other means than by the facts and thoughts lying open or hidden in the Scriptural narrative. Arriving at climaxes, when important direct utterances, especially of the Lord Himself occur, it conforms even to the exact wording of Scripture.

More specifically, the story to be properly told must meet the following conditions: it must be simple, intelligible, intuitional, vivid, and it must have movement. It must be simple in regard to language and tone. Not long but short sentences; simple sentences and not complicated
\quellseite{500}
ones with intertwining clauses. Nor should the catechist cultivate an unnatural and artificial tone, but a natural one, as if a mother tells a story in the family circle. The story must be intelligible; for though it be simple in language and tone, the children may still be unable to understand it, because there is not enough adaptation to the child’s vocabulary, or because words are used which may be familiar as to sound, without, however, producing clear concepts in their minds (p. 439). As a case in point, the teacher should not merely say: “Jehovah blessed the Egyptian’s house for Joseph’s sake”, but should show in detail wherein such blessings consisted: perhaps in this way: While Joseph was overseer, the horses were never hurt; the slaves did not quarrel or steal; and everything sown upon the fields bore fruit, sixty, seventy, and a hundredfold. He must not content himself with repeating the statement of Scripture that Saul was given letters to Damascus; for even pupils of the upper grades do not know what is meant by credentials, which are here meant, and thus they will miss an important feature of the situation. “Our instruction is often so utterly devoid of joy and success because we presuppose something that does not exist” (Schueren, p. 12). How is attention to be aroused; how the interest to be awakened; how the recitation to become acceptable if so much in our narration has remained unintelligible to the children? The narration, furthermore, must be intuitional. “To use the map of Palestine in Biblical History does not make the instruction intuitional. When the child, as the story is told, sees with the eyes of his mind what is heard with the bodily ears; when the external process becomes an internal one for the child, it is then that instruction is intuitional. Brush and pigment are not the only means of painting, nor chisel and hammer of moulding: these things can be done also with words: there is such a thing as a plastic art in language” (Schueren, p. 16); comp. pp. 204 f., 256, 260 f.). To tell a story intuitionally we must pay attention to the way children and mothers tell stories: the Bible, too, contains many sections that can serve as examples. For one thing, it pays attention to the lesser features (Abraham bowed down; let now a little water be fetched and wash your feet; here I am, my son!; Joseph shaved himself and changed his raiment when he left the dungeon; observe the plastic force in Gen. 22!); there are never any omissions (when Abraham entreated God repeatedly in behalf of Sodom, it does not say: He entreated God several times; but the six entreaties are given in accurate detail, and to this day each renewed entreaty causes the child’s sympathy to mount higher); definite numbers are given (five loaves, two fishes, twelve baskets); it does not deal in collective terms, but mentions
\quellseite{501}
concrete objects (e.g., in 1 Sam. 17:17, it does not simply say food, but enumerates the individual articles constituting the food); it never, or at least very seldom (as in Acts 1:1 ff.), deals in generalizations; it does not employ indirect speech with its difficult constructions, but direct speech (not: Jesus said that she should go in peace, but: Go in peace! not, God asked where Adam was, but, God said: Adam, where art thou?). All this enhances intuitional presentation. It is true that the narrative is thus made longer; but that is of no consequence: while breadth is not always an evidence of intuitional presentation, intuitional presentation does require a certain amount of diffuseness. It dare not omit miniature painting, whereby, at the same time, the phantasy of the child is given most acceptable food (pp. 210 f., 251, 256). The outward circumstances should be delineated in detail. Explanatory comments, for instance, such as are to shed light on the state of culture at the time of the event under consideration, should always, as much as possible, be joined to the persons of the story under discussion, unless they have been disposed of in the “preparation”; only thus can they really be made to enhance the intuitional presentation and to arouse the interest of the pupils (p. 257). Under no circumstances should they be accorded recognition as independent factors. While they may be allowed to retard the progress of the action, they must neither interrupt it nor ever crowd the essential features of the story into the background. The ivy twining around the pillar and all the details of decorative art to the right and the left are very fine, but the pillar itself is of greater importance. Special care is required for the illumination of inner processes, that is, the thoughts and sentiments of the actors in the story. The catechist must mentally transplant himself into their position; must envisage their pleasures and their burdens, must visualize their emotions, their fears and hopes now rising, now subsiding, and thus trace the act to its beginnings; he must try to catch a glimpse of the inner conflicts as they end in defeat or victory, and eventually take visible shape in actions; he must take into consideration the consequences of the good as of the evil deed; only thus can he prompt the soul of the child by word, voice, and eye to vibrate in sympathy with the scenes inwardly beheld.—Fankhauser: “When we tell of the campaign undertaken by Abraham for Lot’s rescue, we may fancy how the tempter whispered to him, ‘Serves Lot right! What tricks he used to play on you! Now he has gotten what is coming to him for his selfishness! Let him see how he gets out of the trap!’ But see how Abraham
\quellseite{502}
turned down all such advice and only thought of Lot’s misery—in chains, in fear and danger of death: and he his brother’s son, commended to his care by his dying father. ‘But what will you do, Abraham? With your 318 servants you would fight the four kings with soldiers probably a hundred times as many as yours!’ ‘God goes with me, if I do what pleases Him. If He goes by my side, we shall be the stronger of the two.’ ‘But, Abraham, you have never done anything but to herd flocks, while those men are veteran warriors!’ ‘I trust in God, He will show me what to do’. Zurhellen: ‘Cain’s fratricide, as pictured in the Bible-text, is, for the children, a ghastly and incomprehensible deed. Envy is the motive; and envy should be the subject of the subsequent penetration. But in that case the narration must have succeeded in making clear the nature of envy, how it probably originated, how it gradually corroded the soul, and what terrible results it will have unless opposed by a strong will. Therewith we come upon the second pedagogic motive of great value: ‘Unto thee shall be its desire, but do thou rule over it!’ Self-control is surely one of the most fruitful subjects of religious and moral instruction; and the story before us supplies a welcome opportunity to speak of it to the children. But in that case the presentation should also contain a reference to the conflict of the good will in Cain against sin and make its ultimate defeat conceivable.’ All this is necessary if the narrative is to be intuitional. The catechist will be successful in shedding light upon the outward as well as the inward situation, especially if he observes those inconspicuous hints contained in the Bible text which so often elude the attention of the cursory reader (cf. Gen. 7:21.22; Acts 9:22; above p. 474; or Gen. 15:1 in connection with Genesis. 14). With the instinct of a hunter the catechist must be after those hints; whatever in them appeals to the intuition he must track down; focus his interest upon them, and exploit them. But these processes—the delineation of the external conditions as well as the illumination of the promptings in the inner world of the actors, can enhance the intuitional character of the narrative only when the catechist associates and compares them with such external conditions and experiences of the inner world as the children are already conversant with. This is the germ of truth in these attempts at modernizing the stories of the Bible. With a loud voice they call to the catechist: “Do not forget the connecting threads between the home of the child and the home of your Bible stories; do not forget the fulcrums for the apperceptive faculty (p. 212) in the sphere of the children’s
\quellseite{503}
intellect and experience; look for them; turn them to account if your narrative is to take life in the soul of your child and quicken it to an inner experience of the same.”—If we name, in addition to intuitiveness, also \emph{vividness} as an essential attribute of an effective narration, we have in mind not only a noble, animated delivery, nor merely Fankhauser’s definition: “Vividly to narrate means so to narrate that the story lives in us and we in the story,” but also what Schueren emphasizes when he says that the catechist must have experienced the vital power of the story in his own heart, and that his narrative must become a vivid testimony of the divine life within him.—Finally, narration must have \emph{movement}, it must aim for the goal, it must be \emph{zielstrebend}. This means that it must not tarry too long here or there, least of all with external features; there must be no coloring for the sake of the color; but everything must be viewed in reference to the final goal to be attained, i.e., the precise religious or moral truth in the story which is to be set forth; and everything must be left out, no matter how beautiful, which does not serve the end in view—the realization of that truth.

If the story is told in the manner here set forth, a quickening of the emotions through the intellect will be the result; the story will appear precious to the soul; it will become a factor of its inner experience. The result will be an “ideal association” (p. 260) with the characters of the Bible, a holy \emph{συμπάθεια} (cf. pp. 226 f.). The children think, feel, grieve, pray with them; they plan, strive, suffer, believe, hope with them; they sing, rejoice, shout with them. What a power for culture, education, discipline that means, is known to every pedagogue (pp. 252, 254)!

Last of all, the question remains to be answered whether a story is to be told as a whole or in sections. The latter method still has the most advocates in the present, and the author has not hesitated to give an example of a story told in installments (see practical examples). In case a story is of undue length and makes two or three sharp turns it may be convenient to tell it in sections, and at the end of each section to reduce it to the words of the school text, to have the children reproduce it in their own words, and to let them find the proper heading for the section before proceeding to the next section. But when the story is a
\quellseite{504}
natural catechetical unit and not too long, the effect upon the emotions and the will is not subserved at all when this method is followed; we are convinced that the uninterrupted narrative is every time much surer of a telling effect; and this is decisive for the catechist. In the Biblical History by the author, an outline of every story is given at the conclusion; but this is done altogether in the interest of the pupil, to aid him in reviewing it and committing it to memory. When the whole story has been exhaustively told, the pupils should be directed to open their Biblical History and read the story or recite it extemporaneously, if the latter can be done without torture.—Practical examples are found at the conclusion of the section.

\section{4. Penetration}

With presentation, if effectual, the most important feature in the treatment of a Biblical History has been disposed of. A picture, clear and intuitional, of the story, has been painted for the children from both aspects, external and internal; their soul has been warmed and their will has been set in motion. But the process is as yet far from complete: unless other features are added its effect will be lost too soon and hardly ever cause more than a fleeting mood or sentiment. Instruction and training, however, dare not rest content with the production of religious and moral moods. The object to be attained is: Clear convictions, moral and religious value judgments, a consciousness of the Christian life as a real life—as the receiving of God's grace, as a struggle against sin, as a life of work according to God's will. That this effect, produced by the story upon the child's soul, be perpetuated, clarified, confirmed, deepened, and made fruitful unto life, the children must be directed once more, and that at leisure, to meet the characters of the story face to face; carefully to scrutinize their thoughts and motives—the whole complex of the mental and moral forces impelling them to action; to examine them from the
\quellseite{505}
standpoint of their religious and moral character; to trace carefully the consequences of wicked and of godly conduct; to make the relation of God to man, and of man to God a subject of meditation, so that faith, love, fidelity, unselfishness, etc., may in their majesty confront and captivate them; and that unbelief, unfaithfulness, hatred, envy, selfishness, lying, etc., may stand revealed in their disastrous and abhorrent hideousness. This is the purpose of penetration. It is to gather in leisure all the illuminating and cheering, but also all the repellent and dreadful, rays of the story; to let them exert their inherent force with power, and thus to turn to account the educative elements of the story for the training of the young to the full extent possible. Thus in penetration the last two constitutive elements of each story rather than the first two (p. 473) must be emphasized. For here the catechist comes to speak of the outward progress of the story and its external frame-work only in order that the catechumen may gaze upon the inner world of thoughts and motives, recognize moral and religious values, observe their vital importance, and gain respect and desire for them. When this is understood, it must be clear that penetration deals with persons, especially the leading persons of the story, in accordance with the fact that formative, educative force is mostly found in the personal, biographical element (p. 256). Of whom have we heard in this story? Or, whom do you like best in this story?—these are the questions most naturally asked when the penetration begins. When light is shed upon the thoughts and acts of those persons, the teacher should see to it that the children apply no other standard of judgment than that of God. While sinful acts are by no means excluded from illuminating discussion, care must be taken not to go unduly into detail in the delineation of sin, lest destructive germs are deposited in the children’s hearts. If such guidance to religious and moral value judgments is given by
\quellseite{506}
one who himself loves God and His will, it is by no means, as some maintain, the critical spirit, with its desire to judge others, which is awakened, but rather the children's conscience; that which is good in itself has been made to appear good and precious also to them; and what is evil in itself has been made to appear detestable. While shedding light upon men and their doings, the catechist should never forget shedding light upon God and His doings. His attributes and nature should be brought out in clearest light at the hand of his actions. This is the stage of instruction in which the catechist, in conformity to the example of Christ (John 17:6), is to reveal to his pupils the name of God; that is, he is to withdraw the veil, so that, behind His acts, they see God Himself in His divine nature and recognize Him as absolute personality and holy Love. (Comp. M. Reu, \emph{Unsere Erziehungsarbeit im Lichte von Joh. 17, 6}, Kirchliche Zeitschrift, 1911, pp. 453–466). It is often desirable that the most important points brought out in penetration, for the purpose of impressing them indelibly upon the children's minds, should be pronounced in unison by them and written upon the black-board in the order of their development. When this is done, an easier survey is attained, the final summing up is made easier, and the service of the eye is enlisted in order to secure a more lasting effect for the educative elements found in the story (comp. the complete practical example at the end of this fifth section).

If advisable from the didactic and pedagogic stand-point, there may be joined to penetration what has been called by the disciples of Herbart “\emph{association}” or “\emph{comparison}”; but only in that case. For just the constant—and hence sometimes unavoidably artificial—employment of association and comparison, has, perhaps not without reason, largely occasioned the discreditation of the formal step theory. What is to be done is that the children are made to recall
\quellseite{507}
for themselves the other stories in which the same features in the acts and nature of God, in the doings, experiences, reflections of men, in the relation of cause and effect, have aroused their attention. Whether such stories are taken from Biblical History, Church History, the Reader, or, on occasion, even from other branches; or whether they are taken from the children's own experience, is a matter of no consequence. Such stories being found, the truth developed from the new story is confirmed and enhanced in power; its universal character is recognized by the child, and thus it becomes in his eyes an incontestable certainty, and a factor ever to be reckoned with. It is hardly necessary to quote examples; may it suffice to call attention to the importance of trust in God as shown in the life of Abraham, of Moses, of David; to the effectual prayers of Jacob, of Moses, of Solomon, of Elijah, etc. While such comparison is not always required, there is another feature of discussion, called \emph{generalization}, which is indispensable. Invariably the main truth developed from the story should be summarized in the form of an easily remembered, classical sentence, which may be borrowed from Scripture, the Catechism, or the hymnodic treasure of the Church. Here everything is gathered as in a focus. Where it was possible previously successfully to employ the step of comparison, it is natural that the truth of this summarizing sentence is impressed even more deeply upon the minds of the children. It becomes the focus in which the truth not of one but of many stories is gathered. That such truth dare not be mechanically grafted upon the presented story, but should vitally grow forth from within, follows from what has been said. If presentation always should point “goalward”, penetration should do so even more; every constituent part of penetration should lead in a straight line toward the passage in which ultimately everything is gathered together, so that this summing up is not
\quellseite{508}
only the end but the consummation of penetration. While, in penetration, occasions arrive for the introduction of illuminating Bible verses, stanzas of hymns, or parts of the Catechism, these must yield precedence to the golden text, the ultimate goal. Since the work of memorizing must be done largely in school, especially in the lower grades, this golden text must finally be drilled by recitation in unison.

As to the \emph{outer form} assumed by penetration, not much remains to be said. It cannot be any other than a natural didactic conversation. While the didactic question has already been an occasional feature of the descriptive and detailed narration, the principle of independent activity requires that penetration should proceed entirely in the form of questions and answers, without, of course, excluding the occasional employment of the acroamatic form. The practical example based on the story of Zachaeus (at the conclusion of section V) should be converted into the form of a didactic conversation.

\section{5. Application}

The necessity for an application to follow penetration and, likewise, the aim of the same, has been demonstrated in ch. 32. But what character is the application to bear? At the very outset, it is necessary to caution the catechist against the fallacy that numerous applications are required. In presentation the catechist keeps in view one single aim; and during the progress of penetration he aims at only one summarizing sentence although he may employ a number of Scripture passages or hymn stanzas; in the same manner the catechist should not in the end commit the error of making many applications: the effect of that would be that none has made an impression deep enough to become a fruitful factor in the daily life of the children. What has been said of the preacher, applies to the catechist as well: “There are preachers that take too much for granted; they take a chest full of truths along with them upon the
\quellseite{509}
pulpit, just as someone might take a chest full of nails; and then, on the supposition that the hearers are posts, they take out a nail and expect it somehow to enter the post of itself. But that is not the right way. They must take the nail, hold it to the post, hammer it in, and then clinch it on the other side. Not until then have they the right to expect that the great Master will so fasten the nail that it will not fall out again.” This is true in an even higher degree of the catechist; for children are less able than adults to take home with them a mass of loosely applied truths and to make them factors in their daily life. And just as but one application is to be made, so this one application should not be a general one but one very definite, clearly bounded on every side; for any application necessarily loses force in the degree that it is general. E.g., when the story of Abraham’s departure for Canaan is to find its application, it should not be formulated like this: “Children, you, too, must believe in God,” but: “though God asks difficult things of us, we still should obey Him, confident that He will be with us and bless us.” Finally, the application should be something more than a mere doctrine. If the latter were the desideratum, the application would mark no progress; for the doctrine taught by the story has already been found by penetration and has been gathered into the golden text. What the catechist has to do, is to point out definite incidents in the life of the children and show, in contrast to their real conduct, what it ought to have been. To this end he will go with them into the past or future of their own life, in order to train the conscience for a recognition of sin, and to awaken in them wholesome purposes and resolutions. Moreover, it is their life, not that of others, to which application will have to be made; for in their present life, as children, they have to begin to adapt themselves in all seriousness to the religious and moral truths of Scripture. Instruction
\quellseite{510}
is to be more than a mere aggregate of directions to the children regarding their conduct in the future. Also as adults they will deem it a matter of course in all things to consult the will of God, if, as children, in their relation to parents and teachers, to brothers and sisters, to the neighbor and the neighbor’s children, in the house and in the garden, while at work and while at play, they were accustomed to conform to that will. The skill of the teacher will have its mettle tested in the attempt to restate the great deeds of the Bible characters in the terms of the children’s humble life. He who deems himself too great for that or even looks upon it as a triviality, may make many an application quite good in itself; but the souls of the children he will not touch: and it is his fault when the children fail to locate the channels into which the vital forces quickened in them are to issue; when the impulses awakened in their hearts through presentation and the penetration evaporate without fruitage, and thus already in youth that cleavage is formed between knowledge and life, worship and every-day life, which afterward, in adult life, proves so difficult of removal. Not seldom a short, well-chosen story germane in every way, is best suited to open the eyes of the children for the duty suggested by the biblical narrative with which they have dealt. Of course, the story intended to clinch the application should not be taken from the Bible; it should be one from life, from the present, from Missions; and it should be a juvenile story though this is not an invariable rule; but it must be only one, not a number of stories.

The catechist will find much material useful for presentation in H. Witt, \emph{Die Biblischen Geschichten Alten und Neuen Testaments mit Bibelwort u. freier Zwischenrede anschaulich dargestellt}, 3 vols., ²1883; G. Staebler, \emph{Jesus, der Schoenste unter den Menschenkindern}, Stuttgart, 1905 (however, everything requires re-statement in terms familiar to children); M. R. Unger, \emph{Beitraege zur erbaulichen Behandlung}
\quellseite{511}
\emph{des Religionsunterrichts auf der Unterstufe}, 1912; C. Stuckert, \emph{Jesusgeschichten}, 1910–14; especially in G. Fankhauser, \emph{Der Weg zum Kinde}, 1915 ff., and M. Reu, \emph{How I Tell the Bible Stories to My Children}, ²1926. M. Affolter, \emph{Kommt, Kinder, hoert}, ²1926. \emph{For penetration and application} in J. Nissen, \emph{\emph{Unterredungen ueber die biblischen Geschichten}}, 1888; E. Thraendorf and H. Meltzer, \emph{\emph{Der Religionsunterricht}}, 5 vols., ²1905 ff.; A. Reukauf and E. Heyn, \emph{\emph{Ev. Religionsunterricht}}, vol. iii to ix, ²1906 ff. H. Spanuth, \emph{\emph{Praeparationen fuer den ev. Religionsunterricht}}, 4 parts, 1909 ff.; but preeminently in G. Friedrich (Rienecker), \emph{\emph{Der Herr bleibt ein Koenig in Ewigkeit. Unterredungen ueber die bibl. Gesch. d. N. Testaments}}, 1911, and G. Fr. Rienecker, \emph{\emph{Gottes Wort bleibt in Ewigkeit. Unterredungen ueber Gesch. des A. Test.}}, 1912. \emph{For application} in G. Staebler, \emph{\emph{Sammlung von Erzaehlungen und Grundgedanken f. d. Unterricht in der bibl. Gesch.}}, 1902 f.; J. Heininger, \emph{\emph{Anwendungen bibl. Geschichten}}, 1900; E. Zeller, \emph{\emph{Andachten fuer Kinder}}, ²1914; some material also in G. v. Viebahn, \emph{\emph{Von der Landstrasse des Lebens}}, 1908; and in the \emph{King’s Highway Series}, edited by H. Sneath, G. Hodges, and H. H. Tweedy, 1917.

As far as the \emph{several grades} are concerned in which the Bible stories are to be treated, but little remains to be said. In the lower grades, of course, the whole process of instruction will be simpler and more child-like. But even here the steps named above can be retained; penetration and generalization may be successfully employed in the lower grade, as has been clearly shown by Fankhauser and Unger. It is quite true, however, that it will at times be sufficient in the lower grades to narrate the story in an intuitional manner, and to make the narration issue in the golden text which sums up its main truth. In the upper grade penetration and application are to be cultivated with particular care. If the same stories recur in all three Courses, as is often the case according to the curriculum for parish schools, given in Chapter 31, it need not invariably be the same truth which is placed in the foreground. The individual stories are at times so rich in content that they suggest a variety of inner aims; and often the mental development of the children does
\quellseite{512}
not permit the specification of the ultimate intent of the story until the Third or Upper Course has been reached. The story of the stilling of the tempest may be treated in such a manner that only this truth becomes obvious: “The Saviour saves His own when danger overcomes them”; but also in such a manner that it teaches: “If God be for us, who can be against us?” The pupils in the First or Lower Course understand the former; the second is required for the pupils of the Third or Upper Course. The story of Zachaeus can be treated from the point of view: “Even the publican is not too great a sinner for the Saviour; He enters his house, for the Son of Man is come to seek and to save that which is lost”; but also from this: “If anyone is in Christ, he is a new creature”; and even from the third viewpoint: “To receive forgiveness of sins from Jesus and to right the old wrong, is the way to true happiness.” The story of the marriage at Cana may serve to impress the children that the Saviour is almighty, for He can turn water into wine; but also that “Jesus glories in being a powerful helper even in the small troubles of daily life.” The consideration of the passion of Christ may stop in the First Course at the bodily anguish; in the Intermediate Course stress may be laid upon the inner anguish occasioned by the disloyalty of the disciples and the people's unbelief; while the most difficult feature may be reserved for the Upper Course: the enduring of the divine wrath. The story of Abraham offering Isaac may be made to show that Abraham loved God more than his son; but also that Abraham's faith was staunch and equal to the very greatest trial; and finally, Abraham and Isaac may be held up as types of Christ and God the Father.

“In Quarta the pupils are shown, in the light of Saul's conversion, “how God defeats and hinders every evil counsel and purpose which would not let His kingdom come.” That class, according to the curriculum, studies this story after having learned of the conflict of the Sanhedrin with the Christians, the martyr death of
\quellseite{513}
Stephen, and the dispersion of the Christians. The assurance that Christ is risen from the dead, that He “lives and reigns” is nothing new for them. This assurance, based on what had been taught in Quinta is now merely confirmed. In Tertia where the Life and Teaching of Paul is treated, it is a different truth that should be brought out. There the theme of the story (outer aim) is not: ‘How is the most furious persecutor of the Church conquered by God?’ but: ‘How did Paul come to faith in Jesus Christ?’ And the answer is: ‘God revealed in me His Son’ (Gal. 1:16). ‘He was laid hold on by Christ’ (Phil. 3:12). ‘By the grace of God I am what I am’ (1 Cor. 15:10). ‘By grace are ye saved’ (Eph. 2:8): this is now the saving truth drawn from the very same story. In Quarta the time is too short to bring out this phase, and the pupils would very likely fail to fathom it. And still different will be the statement of aim and the generalization of the saving truth in Upper Secunda. There the subject under consideration is the apostolate of Paul, the divine origin of his Gospel, the history of the conflict between a free Christian faith and the Christian faith as perverted by Pharisaism, or of the transition from Judaistic Christianity to Gentile Christianity; and Bible verses like 2 Cor. 4:6, 1 Cor. 2:9 f. (in connection with 2 Pet. 1:16–21; John 1:14) now summarize in classic words the saving truth of the self-same story.” While these words of Zange (p. 186) presuppose other educational conditions than those prevailing among us, they constitute a notable testimony to the cause to which we stand committed. Let this be noted by our American sticklers for method, who deem it unpsychological and unmethodical to permit the same story to recur in a “Graded Series”, and who take recourse to all kinds of material to avoid this. Zange, a man of pre-eminent training, both as psychologist and as educator, goes so far as to add: “Another proof how necessary it is to return to the same material in the successive grades.” Cf. p. 425.

In the Sunday School, the narration of the story can naturally move toward one aim only because it is intended for all the children in the Second Course (cf. p. 420); but the higher aims can at least be suggested to the pupil by the “Questions for Review” in the Lesson Helps for the Intermediate and Junior Departments when the teacher of the class is not mature enough to lead up to them himself.
\quellseite{514}
The method of teaching here recommended requires much time: in most instances one hour will not suffice for all the successive steps. Nor is it necessary to mount them all in one hour. In the \emph{parish school} the work of the day may close with the presentation and the subsequent reading of the story from the textbook on Biblical History. The pupils are expected to read the story several times at home; and in the next lesson penetration and application take the place of the mechanical questioning which is still widely in vogue. In the \emph{Sunday School} the story is told to the whole Second Course or Department in a detailed manner, as explained above. Because such narration is designed to arouse the feelings of the children and to move their will, it should \emph{never be omitted}; the preparation may be left out (but not the statement of aim) for lack of time. After the story has been told, the several grades may read it under the supervision of the teachers from the graded Lesson Helps, and the Bible verses will be drilled by recitation in unison. On the following Sunday, after the Superintendent has presented and the class has read the new story, the teacher of the individual grade reviews the story of the preceding Sunday with the aid of the review questions appended to every lesson of the \emph{Wartburg Lesson Helps}. Penetration and application of the story taught in the Sunday School may occur in the Saturday School (cf. p. 429). The same method recommended above for a regular parish school can without difficulty be observed in the partial parish school, comprising the children of 12 and 13 years. It is true, not every story can here be thoroughly studied according to our method although a lesson in Bible History is assigned for every forenoon; but that is no great loss when one has done his duty in the Sunday School and Saturday School. Some of the stories are taken up merely by way of review, leaving for thorough treatment only the more important lessons upon
\quellseite{515}
the Sunday School calendar, and those which in Sunday School had been omitted entirely.

When the end of a group of stories has been reached, a review may be undertaken: a detailed sketch of the outstanding character may be drawn, or the principal characteristics of God and the important persons may be assembled, always with a view to the training of souls (p. 483). A sketch germane to the points here made is found among the practical examples at the end of this section.

At times \emph{Bible reading} may be connected with the treatment of the Bible story. In connection with important phases of the discussion, especially when the steps of comparison and generalization have been reached, the upper grade may read—directly from the Bible instead of the textbook—a Psalm or another section that comports well with the truth developed (cf. p. 371; for instance, in connection with the story of creation, Ps. 104 and 8; the passage through the Red Sea, the song of Moses and Miriam; David’s shepherd life, Ps. 23; the story of his fall and subsequent repentance, Ps. 51 and 32; the destruction of Jerusalem, Lamentations; at the turning points in the history of the kings, passages from the prophets; the exile, sections from Isaiah 40–66; the return of Israel from Babylon, Ps. 126; the episode of Cæsarea-Philippi (Matt. 16), Isa. 49; at the foot of the cross, Isa. 53; in connection with the ascension, Ps. 110; on Pentecost, the pertinent section from Joel, etc.).—That the most important facts of \emph{\emph{Biblical Geography, Natural History}}, etc., should be associated in part with the preparation, in part with presentation and penetration, follows readily from the foregoing. There are places where such facts should be brought out with telling effect. When, e.g., the journey through the wilderness is the subject, let the needed geographical information about the peninsula of Sinai be given; when the class enters Canaan with Israel, an accurate survey of the
\quellseite{516}
country should be afforded; when it is taught that God gave Israel every blessing promised, an exhaustive picture of Israel’s life and happiness should be painted (cp. Practical Examples); later the class should follow Israel to Babylon and back; the Christmas story with its mention of Nazareth, Bethlehem, and the land of Judæa furnishes the occasion for a discussion of the division of the Holy Land in that period of its history, which should be illuminated by constant reference to the map; the mention of Emperor Augustus suggests the same treatment of the Roman Empire; when the pupils travel with the twelve-year-old Jesus lad, the journey should be traced upon the map; Paul’s missionary journeys can not possibly be taught without that adjunct. Places like Jerusalem, Bethlehem, Nazareth, and Capernaum lend themselves to geographical illumination (cp. Practical Examples).

Finally, concerning the use of \emph{Biblical pictures} in connection with the several stories, the following remarks are in place. That the use of biblical pictures deserves encouragement as a stimulus for the intuitive faculty, goes without saying; however erroneous the opinion is that the mere use of pictures constitutes intuitive instruction. There are catechists who use pictures diligently, and yet fail to impart intuitional instruction; for the inward intuition, as above described (pp. 478 ff.), is more important than the employment of the finest pictures. This much is true, however, that the internal intuition is enhanced and prolonged when the external intuition is facilitated by a really good picture. But when is the picture to be exhibited and put to use? Before the presentation, during the presentation, or in connection with the application? Each of these three views has found its advocates. The exhibition of the picture before the presentation presupposes that the picture shall be made the starting-point for the whole process of instruction. For the so-called intuitional method of instruction in
\quellseite{517}
the home or lowest grade of the Sunday school that is the proper method; for there the inward intuition is still secondary (compare \emph{Wartburg Lesson Helps}, first course; Wangemann in his \emph{Zwanzig Anschauungsbilder fuer den ersten Unterricht in der Biblischen Geschichte} and in his \emph{Anweisung zum Gebrauch der 20 Anschauungsbilder} goes too much into detail). But in the same degree in which the mental development of the children progresses, facilitating as well as requiring the inward intuition, this method loses its usefulness. While the detailed narration is in progress, the children are still so occupied with the picture that they are not open to the inward intuition intended by the narration; and even the teacher may tend to confine himself to the cultivation of a purely outward intuition.—He who puts the picture to use in connection with the stage of application, is actuated by notions far removed from the purpose of biblical instruction and based upon the presumption that the pupils are possessed of a degree of mental maturity that is ordinarily rare even in the Upper Course of our parish schools. The children are here expected to express an opinion upon the picture and its most salient features, upon the ideas conveyed by the artist upon canvas; this tends to place the picture in the service of esthetic culture.—When we consider that the picture, if it is at all valuable from the catechetical standpoint, represents the climax of the action, it is not likely that any other use of it can be justified than its exhibition during the illuminating narration; that is, just when the climax has been reached and the inward intuition has already taken place. In that case the picture does not divert the attention of the pupils but rather gives strength and depth to the inward intuition, which is now aided by the eye. Having reached the climax, the narrator can properly make a pause and give a little time to a discussion of the principal features of the picture. If the latter crystallizes as it should, the thoughts of the
\quellseite{518}
characters of the story, there will be occasion to recur to it during the penetration.

Among the series of pictures which can be put in the service of the schools, those by Julius Schnorr von Carolsfeld still deserve first place, although this pertains rather to those on the Old Testament. Heilmann, as late as 1911, says: “No one else could bring so much artistic power and religious understanding to bear upon the sacred material as he. In regard to magnificence of composition and passionate interpretation these Bible pictures have hitherto never been surpassed. . . . To illustrate the story of creation, of the patriarchs and prophets, no grander illustrations are to be found. A valuable feature of pedagogic and religious moment is their immaculate moral character and sublime religiousness.” Another merit is that they usually represent the climax of the action. For the New Testament the series by Hoffmann and Schramm excels. There are also good features in the collection of Morgan and Copping; the series by Karl Schmauk (equally realistic, but superior as regards execution of detail), and in the collection made by the Gesellschaft für christliche Kunst in Munich. Recently the Catholic Fugel has produced a good series (twelve Old Testament and as many New Testament pictures have been published to date). Among the pictures illustrative of the parables, those by E. Burnand have justly won renown. Compare the instructive monograph of the Catholic Alfonse Heilmann, \emph{Bibel-Bilder. Gedanken zur religionspädagogischen Wertung biblischer Kunst}, Kempten, 1911.
